\section{Discussion}
\label{sec:discussion}

\subsection{Software architecture as a no-tradeoff reliability lever}
The central empirical result -- a mean relative CV difference of
$0.070\,\%$ between firmware versions, an order of magnitude below typical
replicate variability for this instrument class -- supports a narrow but
useful claim: at least for the scope refactored here, improving firmware
software architecture did not cost measurement correctness. This matters
specifically because the instrument's AFE is commercial, off-the-shelf
silicon~\cite{devices2024ad5940}. In a discrete-op-amp potentiostat design,
a change to the measurement code is entangled with the analog signal chain
in ways that make ``did this change the measurement'' a hardware question
requiring bench characterization of gain, offset, and noise. Here, the AFE's
transimpedance gain, ADC, and DFT engine are fixed, characterized silicon;
what the firmware controls is register sequencing and arithmetic, which
means a refactor that is provably behavior-preserving at the code level
(verbatim extraction, not rewriting) can be expected, and was confirmed, to
be behavior-preserving at the measurement level too. This is the concrete
version of the argument sketched in Section~\ref{sec:intro}: a commercial
AFE converts ``is this instrument reliable'' from a mixed
hardware-and-software question into one where software architecture is the
dominant remaining lever, and that lever is nearly free to pull -- ordinary
refactoring, not a hardware redesign.

\subsection{Relationship to metrological validation}
The companion manuscript~\cite{clement2026linearity} demonstrates a case
where the opposite failure mode occurred: numerically well-behaved code
($R^2 = 0.99999647$ on an Ohm's-law fit) hid a $424\,\%$ calibration
error, because the bug was in a declared constant, not in the code's
control flow or structure. Taken together, the two studies show that
architectural quality and metrological accuracy are separate risk axes.
Duplicated, poorly encapsulated code can compute a correct answer, as it
did here; conversely, cleanly structured code can still compute a wrong
answer if a physical constant is wrong, as the companion manuscript found.
Neither property substitutes for the other, which is why this paper treats
the refactor's validation as a check against regression on a
\emph{previously established} correct baseline, not as a claim that
architecture improvement alone establishes accuracy.

\subsection{Outlook: validation on a real electrochemical solution}
\label{sec:outlook-real-solution}
Every result in this paper uses a passive dummy load -- the Randles
network or the plain resistor -- by design: Section~\ref{sec:protocol}'s
protocol exists specifically to isolate the firmware refactor's effect from
any electrochemistry, since a dummy load's response is predictable from
$R$ and $C$ alone and requires no independent reference measurement to
validate against. That design choice is also this section's boundary: it
means the refactor has not yet been exercised against an actual
Faradaic system, and Figure~\ref{fig:before-after}'s straight CV traces
should not be read as a claim about how the refactored firmware performs
on real chemistry.

\textbf{No solution-phase data was collected for this paper.}
Figure~\ref{fig:dummy-vs-larutan} puts the two side by side, method for
method, so the contrast is a single glance rather than a paragraph: the
top row is this instrument's own real, measured dummy-cell output --
literally the same Randles-cell data as
Figure~\ref{fig:before-after}'s ``after'' traces, replotted here for direct
comparison -- and the bottom row is the idealized, textbook-shape
illustration of what a genuine redox solution's response looks like. The
bottom row is not measured on the instrument, not a rigorous digital
(finite-difference) simulation of the diffusion equation, and carries no
quantitative claim; it exists only so the shape contrast in the top row is
visible without requiring the reader to recall CV theory. Read
column-by-column: the CV columns show the dummy cell's straight line
against the larutan panel's closed-loop, peaked ``duck'' trace (Nernstian
rise to $E_{pc}$/$E_{pa}$, diffusion-controlled tail, $\sim\!57/n$~mV
separation) -- a shape the dummy cell structurally cannot produce, per the
discussion above. The CA columns show the dummy cell settling to a flat,
quantization-limited current against the larutan panel's Cottrell
$i(t)\propto t^{-1/2}$ decay, since only a real concentration gradient
relaxing at the electrode produces that decay. The SWV and DPV columns show
the dummy cell's noise-floor-level, near-flat differential current against
the larutan panels' symmetric bell-shaped peaks at $E_{1/2}$ -- a reasonable
idealization even without a full diffusion treatment, since both are
differential pulse techniques whose subtraction step is specifically
designed to suppress background and tailing, so real pulse-voltammetry
peaks on a reversible couple are, in practice, close to this shape.

The natural next step is to repeat the identical protocol of
Section~\ref{sec:protocol} -- same firmware builds, same four methods, same
before/after and two-condition structure -- against a standard reversible
redox couple (e.g.\ $\mathrm{K_3[Fe(CN)_6]}$/\allowbreak$\mathrm{K_4[Fe(CN)_6]}$ in
supporting electrolyte) and compare the \emph{measured} result to
Figure~\ref{fig:dummy-vs-larutan}'s bottom-row idealization, rather than to
the dummy-cell top row. That experiment would extend this paper's claim
from ``the refactor does not disturb the instrument's response to a
passive load'' to ``the refactor does not disturb the instrument's
response to a real analyte,'' which is the claim ultimately of interest to
a user running this firmware on an actual sample. We flag this explicitly,
and label the bottom row of Figure~\ref{fig:dummy-vs-larutan} explicitly,
so that a reader comparing this paper against one that does include
solution-phase data does not mistake the illustration for a measurement.

\begin{figure}[h]
\centering
\includegraphics[width=\linewidth]{fig3_dummy_vs_larutan}
\caption{HunStat2 read-out, all four methods: dummy cell (top row,
\textbf{real, measured} data, same instrument and firmware as
Figure~\ref{fig:before-after}) vs.\ larutan / real solution (bottom row,
\textbf{idealized illustration, not measured, not simulated}), with the
key diagnostic feature of each method marked on the larutan panel. CV:
straight line (dummy cell) vs.\ closed-loop Nernstian ``duck'' peak
($E_{pc}$/$E_{pa}$ marked, $\Delta E_p\approx 57/n$~mV). CA:
quantization-limited flat trace vs.\ Cottrell $t^{-1/2}$ decay from the
potential step at $t{=}0$. SWV/DPV: noise-floor differential current vs.\
symmetric peak at $E_p=E_{1/2}$ with peak height $\Delta I_p$ marked.}
\label{fig:dummy-vs-larutan}
\end{figure}

\subsection{Limitations}
The refactor was deliberately scoped to the duplication with the clearest,
lowest-risk fix -- the DC-excitation methods (CA, DPV, SWV) and the two
confirmed-dead subsystems. It does not address the audit's other findings:
\texttt{C\_DataStorage} remains a fully public-field parameter store with no
validation, the EIS and OCP classes were not brought into the same
base-class structure (their register sequences differ enough from the
DC-excitation path that doing so safely would need its own validation pass),
and the two-hardware-abstraction-layer question -- separating register
access from measurement logic entirely -- is unaddressed. These are the
natural next steps, not omissions the results depend on. The equivalence
claim itself is bounded by the instrument's own noise floor: SWV and DPV
report differential currents near $10^{-7}$~A, where point-to-point
variation from ordinary measurement noise is a larger fraction of the
signal than for CV, so the before/after comparison for those two methods is
qualitative (same shape, same order of magnitude, no directional bias)
rather than quantified to the same precision as CV. The second-load
corroboration in Section~\ref{sec:results} used one additional resistor
value on the refactored firmware only, not a full before/after repeat at
that load; it is offered as evidence the refactor generalizes, not as an
independent replicate of the primary comparison. Finally, all measurements
were taken on a single physical unit; run-to-run and unit-to-unit
variability at the precision needed to distinguish a $0.07\,\%$ software
effect from unit-to-unit hardware variability is outside this paper's scope
and is the kind of question the companion manuscript's protocol is designed
to answer.
